El plan relaciona los 16 indicadores del capítulo 3 con tareas de diseño y evidencia de cierre. Los responsables son roles propuestos, no contrataciones confirmadas. La definición de una especificación no equivale a su cumplimiento construido.
\clearpage
\subsection{Energía y calidad ambiental}
\begin{table}[H]\centering
\caption{Acciones de implementación y evidencia requerida: energía y calidad ambiental}
{\fontsize{9}{11}\selectfont\singlespacing\renewcommand{\arraystretch}{1.3}
\begin{tabular}{>{\raggedright\arraybackslash}p{1.1cm}>{\raggedright\arraybackslash}p{6.1cm}>{\raggedright\arraybackslash}p{6.1cm}}\toprule
ID & Acción y responsable propuesto & Evidencia de cierre \\ \midrule
E01 & Comparar CB y C3r02 con entradas comunes y nuevas cargas de diseño. Simulación. & Series horarias conciliadas y reducción anual. \\
E02 & Construir referencia EDGE y contabilizar todos los usos y sistemas reales. Especialista energético. & Ahorro operacional mínimo del 20 \,\% frente a referencia EDGE. \\
E03 & Resolver aislamiento continuo, puentes térmicos y humedad. Confirmar clima y compacidad. Arquitectura. & Memoria de U del conjunto y detalles constructivos. \\
E04 & Comprobar cada orientación y ajustar aleros o lamas sin anular luz natural. Arquitectura. & SHGC modificado $<$0,20 en al menos 90 \,\% de superficies aplicables. \\
E05 & Especificar LED y verificar regulación, lux y cobertura del control. Diseño eléctrico. & Potencia $\leq$8 W/m\textsuperscript{2}; fotometría y control según matriz. \\
C01 & Evaluar horas ocupadas por zona y separar funcionamiento mecánico y natural. Simulación. & PMV o método adaptativo aplicable; porcentaje espacial y temporal. \\
C02 & Dimensionar aire exterior, filtración y recuperación; verificar condiciones de excepción. Diseño HVAC. & Memoria de caudales y equipos según T4 y normativa aplicable. \\
C03 & Calcular sDA y ASE con la misma malla, ocupación y plano de trabajo. Simulación lumínica. & Resultados anuales y tratamiento del deslumbramiento. \\
\bottomrule\end{tabular}}
\fuente{Elaboración propia con los umbrales y condiciones de la matriz del capítulo 3.}
\end{table}
\clearpage
\subsection{Recursos, sistemas y operación}
\begin{table}[H]\centering
\caption{Acciones de implementación y evidencia requerida: recursos, sistemas y operación}
{\fontsize{9}{11}\selectfont\singlespacing\renewcommand{\arraystretch}{1.3}
\begin{tabular}{>{\raggedright\arraybackslash}p{1.1cm}>{\raggedright\arraybackslash}p{6.1cm}>{\raggedright\arraybackslash}p{6.1cm}}\toprule
ID & Acción y responsable propuesto & Evidencia de cierre \\ \midrule
A01 & Seleccionar griferías y sanitarios eficientes con caudales a presión de ensayo. Diseño sanitario. & Balance de usuarios y ahorro $\geq$20 \,\% respecto a referencia EDGE. \\
M01 & Medir cantidades y comparar EPD con igual unidad funcional y límites. Arquitectura y presupuesto. & Balance EDGE de materiales; reducción $\geq$20 \,\% demostrada. \\
M02 & Solicitar ensayos de emisiones para acabados, adhesivos y aislantes. Especificación y compras. & Expediente de todos los productos de la intervención. \\
E06 & Determinar SRE, área útil sin sombras y demanda eléctrica. Diseño eléctrico. & Potencia mínima según T2; diseño eléctrico y estructural. \\
E07 & Seleccionar equipos con SEER verificable y revisar dimensionamiento. Diseño HVAC. & SEER $\geq$6 y sobredimensionamiento $\leq$15 \,\%, según aplicación. \\
G01 & Asignar medición de energía, agua e indicadores interiores. Operación. & Plan con responsable, frecuencia y acciones correctivas. \\
G02 & Comparar suministro, transporte, salinidad, mantenimiento y reposición. Presupuesto y operación. & Cotizaciones y evaluación de ciclo de vida y logística. \\
P01 & Comprobar área libre respecto al suelo, seguridad y protección frente a lluvia. Arquitectura. & Abertura efectiva $\geq$4 \,\% en oficinas cuando aplique A6. \\
\bottomrule\end{tabular}}
\fuente{Elaboración propia con los umbrales y condiciones de la matriz del capítulo 3.}
\end{table}
\clearpage
\subsection{Lineamientos transferibles al contexto insular}
La secuencia replicable consiste en reducir primero las cargas evitables, comprobar después la respuesta de cada espacio y seleccionar finalmente sistemas y materiales con evidencia de suministro y mantenimiento. La potencia de iluminación debe vincularse con la tarea visual; la protección solar, con la orientación y la luz natural; y la ventilación, con el clima y los caudales efectivos.

Las transmitancias o fracciones de apertura de este ensayo no se trasladan automáticamente a otros edificios. Cada aplicación requiere geometría, clima, ocupación y condiciones constructivas propias. En Galápagos, la reposición, el transporte y la exposición ambiental forman parte de la selección técnica y de su viabilidad económica.
